\section{Summaries of the three reference papers}\label{app:summaries}
\label{sec:summaries}

These summaries are in our own words, for internal reference. They are
deliberately compact, roughly one page each, and emphasise the features that
matter for our comparison in Section~\ref{sec:lit}.

\subsection{\citet{ARS2024}, ``Managing an Energy Shock''}

\citet{ARS2024} ask how an energy-importing economy should respond to an energy
price shock, and whether household heterogeneity changes the answer. The model
is a small open economy HANK built on the exchange-rate model of Auclert,
Rognlie, Souchier and Straub (2021), extended with an energy good. There are
three goods: a domestically produced home good, an imported foreign good, and
imported energy. Households face idiosyncratic productivity risk, cannot insure
it, and hold a single asset through a mutual fund; preferences are CRRA in a
CES consumption basket that nests energy against a home/foreign bundle, with a
low energy substitution elasticity ($\eta_E=0.1$) and a total substitution
index $\chi=0.3$. Crucially, all households share the \emph{same} basket and
the \emph{same} price index: the energy technology is fixed and homothetic.
Wages are sticky with a real-wage-stabilisation motive; the monetary baseline
is a rule that holds the real interest rate constant. Energy is supplied by
price-taking foreign firms with an intertemporal extraction margin, so in the
single-country equilibrium the world energy price is exogenous.

The analytical core is a decomposition of the output response into an
expenditure-switching channel (positive) and a real-income channel (negative)
scaled by the matrix of intertemporal MPCs. In the complete-markets
representative-agent benchmark only expenditure switching survives, so the
shock is expansionary; with realistic MPCs and low $\chi$ the real-income
channel dominates and the shock is contractionary, their central result, and
the one we inherit. A neutrality result holds at $\chi=1$, where the two
channels cancel. With real wage rigidity a wage--price spiral can appear, but
the real wage still falls.

On policy, AMRS study monetary tightening and three fiscal tools, energy
subsidies, targeted transfers (indexed to counterfactual energy use), and
untargeted transfers, all deficit-financed and repaid by a labour tax. For a
country acting alone, subsidies are the most effective domestic stabiliser and
tame inflation (they cap the price that drives the real-wage loss); transfers
support demand but are more inflationary. The decisive twist is
\emph{cross-country}: monetary tightening has a positive externality (it lowers
world energy demand and price), while subsidies have a large negative
externality (coordinated subsidies make world energy demand nearly inelastic,
spike the world price, and become self-defeating). Their policy recommendation
is coordinated tightening plus transfers targeted to the poor, avoiding
subsidies. They also show the impulse responses are not strongly nonlinear at
empirically relevant shock sizes.

For our purposes the key features are: fixed homothetic energy technology, a
single energy price and CPI, the constant-real-rate monetary rule, and the fact
that the only ``cost'' of a subsidy in their world is a cross-country
externality, there is no transition margin.

\subsection{\citet{Bayer2026}, ``Hicks in HANK'' / ``Redistribution Within and
Across Borders''}

\citet{Bayer2026} evaluate the fiscal response to the European energy crisis in
a two-country HANK model of a monetary union, calibrated with Germany as
``Home'' (one-third of the union) and the rest of the area as ``Foreign''
(calibrated to Italy). The model is a medium-scale, two-asset HANK: households
hold a liquid union-wide bond and illiquid capital, face idiosyncratic income
risk, and have GHH preferences; there is an investment margin with capacity
utilisation. Energy enters both production (a CES bundle with the
capital--labour aggregate, elasticity 0.2) and household consumption (a CES
bundle with the physical good, elasticity 0.1). A distinctive and, for them,
central feature is within-country heterogeneity in the \emph{energy share}:
households are of high- or low-energy-intensity type, following a persistent
Markov process correlated with income, so that poorer households spend a larger
share on energy and suffer larger real-income losses. Monetary policy is a
standard union-wide Taylor rule; each national fiscal authority taxes labour
and profits and stabilises its debt.

The defining assumption is that the total union energy supply is
\emph{inelastic}: import capacity is fixed in the short run, so a fixed quantity
is set exogenously and a common price clears the market. The crisis is a 20\%
cut in supply for six quarters, in response to which the energy price rises
roughly fivefold, inflation rises about four points, and output falls around
one percent, with consumption falling most for poor, high-energy households.

They compare two policies, each implemented in Home only so that spillovers can
be measured: a subsidy that caps the retail energy price at its pre-crisis
level, and a Hicks/Slutsky transfer indexed to pre-crisis energy consumption
(to households and firms). The subsidy is the more effective domestic
stabiliser but is fiscally expensive and \emph{beggar-thy-neighbour}: by
holding Home's energy demand up it pushes the common price higher, amplifying
the crisis in Foreign. The transfer stabilises less at Home but generates
essentially no spillover and is cheaper. On welfare (consumption-equivalent
variation), the subsidy \emph{reduces} Home welfare, subsidising energy under a
supply contraction is inefficient because it suppresses both intratemporal and
intertemporal substitution, while the transfer \emph{raises} Home welfare, its
cost being mostly the distortion of the taxes that finance it. Their headline
contrast with AMRS (subsidy more stabilising than transfer for output) they
attribute to the Taylor rule and the two-country structure, versus AMRS's
constant-real-rate rule.

For our purposes the key features are: fixed union-wide quantity (our inelastic
closure), the two-country spillover mechanism that we cannot represent, the
within-country energy-share heterogeneity that we do not have, and the
substitution-preserving property of the transfer, which we share.

\subsection{\citet{Langot2026}, ``The Macroeconomic and Redistributive Effects
of the French Tariff Shield''}

\citet{Langot2026} evaluate the French \emph{bouclier tarifaire}, the 2022--23
cap on gas and electricity prices, in a HANK model of France as a single
economy within the euro area. The model extends Auclert, Rognlie, Souchier and
Straub with energy in both household consumption and production. Households face
idiosyncratic productivity risk and hold a single asset; preferences are CRRA
with a CES energy nest, and the energy substitution elasticity is $\eta_E=0.5$,
markedly higher than in AMRS or Bayer. A distinctive feature is an
\emph{incompressible} (subsistence) level of energy consumption: only energy
above the subsistence level yields utility. This generates non-homotheticity
without non-homothetic preferences, poorer households devote a larger share to
incompressible energy and have a lower price elasticity of energy demand, which
is the source of the distributional results. The ECB Taylor rule is
down-weighted to reflect France's small weight in euro-area inflation. In the
policy simulations the energy price is taken as exogenous (from the
government's forecasts), so France is effectively a price taker facing
perfectly elastic supply.

The methodology is distinctive and worth flagging: rather than clean MIT-shock
impulse responses, the paper performs an \emph{ex-ante, conditional-forecast}
evaluation. Using the sequence-space Jacobian, the authors back out the
sequence of structural shocks (energy price, government spending, transfers,
plus latent demand, markup and fiscal shocks) that makes the model reproduce
the French government's official published forecasts for output, inflation and
the debt ratio through 2027. Policies are then evaluated as counterfactuals
that hold this shock sequence fixed and change one instrument.

The main finding is that the tariff shield supports growth (about 1.9\% average
over 2022--23 against 1.0\% without it) and curbs inflation (5.6\% against
7.2\%), at a fiscal cost near 2\% of GDP per year and a modest rise in the debt
ratio, while containing, though not eliminating, the rise in consumption
inequality even though it is untargeted. The mechanism behind its effectiveness
is a \emph{labour-cost} channel: by holding down the retail energy price the
shield prevents the price--wage spiral, contains the cost of labour, and
thereby supports employment and output. On this macro criterion the shield
beats both a faster wage indexation (more inflationary, worse for output) and a
targeted transfer to finance incompressible energy consumption; the targeted
transfer, however, reduces inequality by more.

For our purposes the key features are: exogenous/elastic energy price (our
price-taking closure), the labour-cost channel behind the pro-shield result
(which our wage block only partly captures), the higher intensive substitution
elasticity, and again a fixed energy technology with no transition margin.
